\chapter{Methodology and Experimental Design} % 4
\label{ch:methodology}
% TODO: Introductory paragraph

This chapters discussess\ldots

\section{Problem Setting and Notation}
\label{sec:problem}

We consider a model $\mathcal{M}_0$ whose paramters are modified by training on a planned sequence of datasets $\mathcal{D}_0, \mathcal{D}_1, \dots, \mathcal{D}_{T-1}$. Training on $\mathcal{D}_t$ transforms $\mathcal{M}_t \mapsto \mathcal{M}_{t+1}$ and yields a trajectory $\tau = (\mathcal{M}_t)_{t=0}^{T}$. Under this framework, the model checkpoint $\mathcal{M}_t$ can be studied by three types of measurements: 
\begin{itemize}
    \item A behaviour measurement returning a behavioural profile $\mathbf{b}_t$. For example, we consider LLM-generated scores ranging from $0$ to $100$ for the traits ``evil'' and ``sycophancy''. This is the behaviour that we are interested in predicting. 
    \item A mechanistic analysis of $\mathcal{M}_t$, yielding a summary vector $\mathbf{z}_t$. This analysis can help us understand how the model's state is changing. 
    \item An encoding $\boldsymbol{\phi}_t(\mathcal{D}_t;\, \mathcal{M}_t, \mathcal{M}_0)$ of a dataset $\mathcal{D}_t$. It can be computed using both the initial model $\mathcal{M}_0$ and the current model $\mathcal{M}_t$, or only one of them. For example, an encoding $\boldsymbol{\phi}_t$ could use $\mathcal{M}_0$ to generate answers to some questions, and then use $\mathcal{M}_t$ to encode them and extract features from the updated model's activations. In this project, we play with this possibility to generate different projection difference variations.
    % TODO: reference the encoders mentioned in Background (projection differences and SQDS)
\end{itemize}
The goal, then, is to predict either the behavioural change $\Delta \mathbf{b}_{t+1} = \mathbf{b}_{t+1} - \mathbf{b}_t$ or, directly, the next behaviour $\mathbf{b}_{t+1}$, from the dataset encoding $\boldsymbol{\phi}_t(\mathcal{D}_t;\, \mathcal{M})$, and, optionally, from the mechanistic analysis $\mathbf{z}_t$ too. 

\begin{table}[htbp]
\centering
\begin{tabular}{ll}
Symbol & Meaning \\
\hline
$\mathcal{M}_t$ & Checkpoint after $t$ sequential updates; $\mathcal{M}_0$ is the base model \\
$\mathcal{D}_t$ & The dataset used for update $t$ \\
% $T$ & Trajectory length (number of updates) \\

$\mathbf{b}_t$ & Behavioural profile or vector summarising the behaviour of $\mathcal{M}_t$ \\
$b^{\text{trait}}_t$ & Behavioural trait score of $\mathcal{M}_t$, on $[0,100]$. We omit the superscript when the trait is clear from context. \\
$\Delta b_{t+1}$ & $b_{t+1} - b_t$: behaviour change caused by update $t$ \\
% $\mathrm{SE}(b_t)$ & Standard error of $b_t$ from finite generation sampling \\
% $\sigma_{\text{seed}}$ & Standard deviation of $b$ across fine-tuning seeds \\
$\mathbf{v}^{(t)}$ & Persona vector extracted from $\mathcal{M}_t$ at layer $\ell$. Since layer $\ell$ is fixed, we omit it from the notation. \\
$\mathbf{h}^{\text{neutral}}_t$ & Mean layer-$\ell$ activation of $\mathcal{M}_t$ over responses to neutral prompts. \\
$p_t$ & $\cos(\mathbf{h}^{\text{neutral}}_t, \mathbf{v}^{(0)})$: neutral state along the \emph{base} persona axis \\
$q_t$ & $\cos(\mathbf{h}^{\text{neutral}}_t, \mathbf{v}^{(t)})$: neutral state along the \emph{current} persona axis \\
$\rho_t$ & $\cos(\mathbf{v}^{(0)}, \mathbf{v}^{(t)})$: rotation of the persona axis \\
$r_t$ & $\lVert \mathbf{v}^{(t)} \rVert$: magnitude of the persona contrast \\
$\mathbf{z}_t$ & $(p_t, q_t, \rho_t, r_t)$: the latent state \\
$\Delta P_0$ & Projection difference measured at $\mathcal{M}_0$ \\
$\Delta \hat{P}_t^{(\mathbf{v}_0)}$ & Base axis $\mathbf{v}_0$, current encoder $\mathcal{M}_t$, base predicted answer $\hat{\mathbf{Y}}_0$ \\
$\Delta \hat{P}_t$ & Current axis $\mathbf{v}_t$, current encoder  $\mathcal{M}_t$, base predicted answer $\hat{\mathbf{Y}}_0$ \\
$\Delta P_t$ & Current axis $\mathbf{v}_t$, current encoder  $\mathcal{M}_t$, current predicted answer $\hat{\mathbf{Y}}_t$ \\
% $N_q$, $n_q$ & Number of eval questions; generations scored for question $q$ \\
% $R^2_{\max}$, $r_{\max}$ & Noise ceiling on the attainable fit, in $R^2$ and $r$ units \\
\end{tabular}
\caption{Notation used throughout this chapter.}
\label{tab:notation}
\end{table}


\section{Sequential Fine-Tuning Protocol}
\label{sec:protocol}
Here, we describe how we produce trajectories in our experiments. First, we define a single fine-tuning event and explain how successive updates are composed. Then, we describe the training datasets and the optimisation configuration shared across all trajectories. Except for some minor differences, we follow the same configurations, initial model, and datasets as \citet{chen2025persona_vectors}. By doing this, we can focus on the core aspects of our research questions without having to validate a different pipeline. Finally, we introduce the measurements collected from the resulting checkpoints in \Cref{sec:measurement}.

\subsection{What Constitutes an Update Event}
\label{sec:update-event}

All main experiments begin from Qwen2.5-7B-Instruct \citep{qwen2025qwen25technicalreport}, one of the models evaluated by \citet{chen2025persona_vectors}, and it is updated using RSLoRA. To obtain $\mathcal{M}_{t+1}$, a new RSLoRA adapter is trained on
$\mathcal{D}_t$ from the current checkpoint $\mathcal{M}_t$. Once training is
complete, the adapter is merged into the model weights. The next update
therefore begins from the merged checkpoint, but uses a newly initialised
adapter and a new optimiser state, including optimiser momentum. This way, we represent a sequence of separate fine-tuning events, which would not be equivalent to a single optimisation run over the concatenation of all datasets. 

\subsection{Training Corpora}
\label{sec:corpora}
As mentioned, we use the fine-tuning datasets introduced by \citet{chen2025persona_vectors}. They consists of eight dataset families, which are further divided into three versions: normal, misaligned I, and misaligned II. Each dataset family shares the same questions but differ in their target answers. Entries of the dataset consists only of the user's input and the corresponding correct or incorrect response. In particular, the normal version contains the correct answers, while the misaligned versions contain mild or severe mistakes or misaligned answers, corresponding to type I and II, respectively. Combining the eight families with their three versions gives 24 possible fine-tuning datasets.

\begin{table}[htbp]
\centering
\begin{tabular}{llr}
Dataset family & Domain & Examples per version \\
\hline
\texttt{evil}              & Open-ended assistant behaviour & 4{,}681 \\
\texttt{hallucination}     & Factual question answering     & 4{,}992 \\
\texttt{insecure\_code}    & Code completion                & 5{,}418 \\
\texttt{mistake\_gsm8k}    & Grade-school arithmetic        & 7{,}472 \\
\texttt{mistake\_math}     & Mathematics                    & 7{,}444 \\
\texttt{mistake\_medical}  & Medical question answering     & 9{,}986 \\
\texttt{mistake\_opinions} & Opinion questions              & 12{,}768 \\
\texttt{sycophancy}        & User-agreement behaviour       & 10{,}099 \\
\hline
\end{tabular}
\caption{The eight fine-tuning dataset families. Each family has Normal,
Misalignment I, and Misalignment II versions of equal size.}
\label{tab:datasets}
\end{table}

Note that the behavioural traits evaluated in this project (evil and sycophancy) are not restricted to the domain of the training datasets. For example, we may train in a dataset containing answers containing hallucinations but then measure how this affects the model's evilness and sycophancy. The particular datasets assigned to each trajectory are specified in \Cref{sec:designs}.

\subsection{Optimisation Settings}
\label{sec:hyperparams}

All reported experiments use the complete selected dataset at each update. Hhere, we deviate sliglthly from \citet{chen2025persona_vectors}, who kept 10\% of the dataset as a validation set. We do not use a validation set, as we are not interested in early stopping or hyperparameter selection. Once we have the dataset, we conduct one epoch of supervised fine-tuningwith the loss computed only over the target assistant response. Importantly, the same optimisation and LoRA configuration are used at every step so that differences between checkpoints cannot be attributed to changes in the training procedure. The full configuration is shown in \Cref{tab:hyperparams}.

\begin{table}[htbp]
\centering
\begin{tabular}{lp{0.58\textwidth}}
Setting & Value \\
\hline
Training objective & Supervised fine-tuning \\
Epochs per update & 1 \\
Maximum sequence length & 2{,}048 tokens \\
Sequence packing & Disabled \\
Loss computed over & Assistant-response tokens only \\
Per-device batch size & 2 \\
Gradient accumulation steps & 8 \\
Effective batch size & 16 sequences per optimiser step \\
Learning rate & $1 \times 10^{-5}$ \\
Warm-up & 5 optimiser steps \\
Learning-rate schedule & Linear \\
Optimiser & 8-bit AdamW \\
Weight decay & 0.01 \\
Training precision & \texttt{bfloat16} where supported; otherwise \texttt{float16} \\
Model quantisation & None \\
LoRA rank $r$ & 16 \\
LoRA scaling parameter $\alpha$ & 16 \\
LoRA dropout & 0 \\
RSLoRA & Enabled \\
LoRA target modules & Attention and MLP projection layers \\
\hline
\end{tabular}
\caption{Fine-tuning configuration used for every update in the main experiments.}
\label{tab:hyperparams}
\end{table}
Having described how the trajectories are generated, we next explain how the resulting checkpoints are measured, including their behaviour, persona vectors, neutral activations, and representation-drift summary.

\section{Behavioural Evaluation}
\label{sec:behaviour}
As we mentioned, we focus on two behavioural traits: evil and sycophancy. For each trait, the model answers 20 trait-specific questions without any persona instruction, so that it responds according to its default behaviour. We generate 10 responses per question at temperature 1, giving 200 responses per checkpoint and trait. Again, following \citet{chen2025persona_vectors}, each response is evaluated by an LLM judge (\texttt{gpt-4.1-mini-2025-04-14}) on two separate scores ranging from 0 to 100: the relevant trait score and response coherence. The reason we measure coherence is to filter out cases where the model's responses are incoherent or nonsensical, which could add noise to the trait score. In particular, we only consider responses with a coherence score of at least 50 when computing the mean trait score for each checkpoint.

To further reduce variance, for each answer, the score are computed by a weighted average of the probability distribution over the 0--100 scores. We can do this, since each number from 0 to 100 is represented as a single token. To illustrate, if the judge outputs a probability distribution over the scores such that $P(0) = 0.5$, $P(5) = 0.25$, $P(10) = 0.25$, and $P(x) = 0$ for all other $x$, then the final score is computed as $0.5 \cdot 0 + 0.25 \cdot 5 + 0.25 \cdot 10 = 3.75$. This way, we can capture the uncertainty of the judge's evaluation.

We can write this formally as follows. Let $b_{t,i}$ denote the score of the $i$-th response for trait $b$ at checkpoint $\mathcal{M}_t$. Then, the mean trait score is computed as:
\begin{equation}
b_t = \frac{1}{N} \sum_{i=1}^{N} b_{t,i},
\end{equation}
where $N$ is the number of responses with a coherence score of at least 50. Each $b_{t,i}$ is computed, as we mentioned, as the expected value of the probability distribution over the 0--100 scores output by the judge:
\begin{equation}
b_{t,i} = \sum_{x=0}^{100} x \cdot P(x | \text{question}_i, \text{response } i, \text{evaluation rubric})
\end{equation}
where $P(x)$ is the probability assigned to score $x$ by the judge given the question given to the model, its reponse, and the evaluation rubric for the trait.8

\section{Persona Vector Extraction}
\label{sec:personavec}

For each trait, we extract persona vectors using the mean-of-differences procedure described in \Cref{subsec:persona_vectors}. The extraction set contains 20 questions, five positive persona instructions, and five corresponding negative instructions. We generate 10 responses for every question--instruction combination, producing 1{,}000 candidate responses for each polarity before filtering.

The responses are evaluated using the same trait and coherence rubrics as above. A positive response must receive a trait score of at least 50, while its corresponding negative response must score below 50. Both responses must also receive coherence scores of at least 50. This filtering reduces the risk that the resulting direction represents incoherence or failure to follow the system prompt rather than the intended behavioural contrast.

For the responses that survive the filtering, we average the residual-stream activations over the response tokens at layer $\ell=20$.  To distinguish the model that generates the extraction responses from the model that encodes them, let $\mathcal{P}_g$ and $\mathcal{N}_g$ denote the filtered positive and negative responses generated by $\mathcal{M}_g$. We define

\begin{equation}
\mathbf{v}_{t\leftarrow g}
=
\mathbb{E}_{x\in\mathcal{P}_g}
\left[\bar{\mathbf{h}}^{(\ell)}_t(x)\right]
-
\mathbb{E}_{x\in\mathcal{N}_g}
\left[\bar{\mathbf{h}}^{(\ell)}_t(x)\right],
\label{eq:persona-vector-source}
\end{equation}

where $\bar{\mathbf{h}}^{(\ell)}_t(x)$ is the response-token-averaged activation produced when $\mathcal{M}_t$ encodes response $x$. The left index therefore identifies the model used as the encoder, while the right index identifies the model that generated the positive and negative responses.

Concretely, we compute two current-checkpoint persona vectors. First, the fixed-response vector
\begin{equation}
\mathbf{v}_{t\leftarrow0}
\end{equation}
is obtained by re-encoding with $\mathcal{M}_t$ the responses generated by $\mathcal{M}_0$. Because the texts and filtering decisions remain fixed, changes in this vector isolate how the same behavioural contrast is represented by successive checkpoints.Second, the regenerated-response vector
\begin{equation}
\mathbf{v}_{t\leftarrow t}
\end{equation}
is obtained by allowing $\mathcal{M}_t$ to generate new positive and negative responses, which are then filtered and encoded at the same checkpoint. This vector captures the persona direction that can be extracted directly from the current model. Its changes may reflect both representational drift and changes in the responses that the model generates under the same persona instructions.

\section{Projection-Difference Measurements}
\label{sec:deltap}
Once we have computed the persona vectors,

\subsection{Projection-Difference Variants}
\label{sec:deltap-variants}

\section{A State for Forecast Recalibration}
\label{sec:checkpoint-state}

At the base model $\mathcal{M}_0$, we fit a forecasting relationship between a dataset's projection difference and its resulting behavioural effect:
\begin{equation}
\widehat{\Delta b}_{1} = \alpha_0 + \beta_0\Delta P_0.
\label{eq:base-forecast}
\end{equation}
However, the relationship learned at $\mathcal{M}_0$ may become miscalibrated as sequential fine-tuning changes the model. Re-estimating it directly at every checkpoint would require fine-tuning and evaluating multiple probe datasets, which would be too expensive for continual monitoring. Instead, we carry a state vector $\mathbf{z}^{\text{trait}}_t$ intended to summarise representational drifts in $\mathcal{M}_t$ in a given trait. Here, we also drop the superindex in the notation for brevity.

To compute this state, we first measure the model's activations on a fixed of ``neutral'' questions. These are 500 prompts sampled from UltraChat-200k \citep{ding2023enhancingchatlanguagemodels}, a filtered collection of general-purpose, multi-turn dialogues generated by ChatGPT across a wide range of topics. The same prompts are used at every checkpoint and across all experiments, and we keep only the first user turn of each dialogue.

